\documentclass[12pt]{exam}
\usepackage[T1]{fontenc}
\usepackage[utf8]{inputenc}
\IfFileExists{lithuanian.ldf}{\usepackage[lithuanian]{babel}}{\usepackage{babel}\babelprovide[main,import]{lithuanian}}
\usepackage{lmodern,graphicx,xcolor,tikz}
\usepackage[a4paper,left=1cm,right=1cm,top=2.5cm,bottom=1cm,footskip=12pt]{geometry}
\usepackage{amsmath,amssymb,tasks}
\newcommand{\irasolaukas}[1]{\begingroup\setlength{\fboxrule}{0.3pt}\setlength{\fboxsep}{3pt}\fcolorbox{black!35}{white}{\rule{0pt}{0.45cm}\hspace{#1}}\endgroup}
\newcommand{\sprendimas}[1]{\par\smallskip\noindent\textit{Sprendimas:}\par\nobreak\vspace{0.1cm}\noindent\begin{tikzpicture}\path[use as bounding box](0,0) rectangle (\linewidth,{-0.7*#1});\foreach \i in {1,...,#1}{\draw[draw=black!25,line width=0.3pt](0,{-0.7*\i}) -- (\linewidth,{-0.7*\i});}\end{tikzpicture}\par\smallskip}
\newcommand{\atsakymas}[1][9]{\par\smallskip\noindent\textit{Atsakymas:}\framebox[#1cm]{\rule{0pt}{0.5cm}}\par}
\pagestyle{headandfoot}\runningheadrule
\pointsdroppedatright\bracketedpoints\setlength{\rightpointsmargin}{1.6cm}
\pointpoints{taškas}{taškai}
\makeatletter
\renewcommand{\points}{%
  \begingroup
    \ifnum\numexpr\@points-100*(\@points/100)\relax>9
      \ifnum\numexpr\@points-100*(\@points/100)\relax<20
        taškų%
      \else
        \ifnum\numexpr\@points-10*(\@points/10)\relax=1
          taškas%
        \else
          \ifnum\numexpr\@points-10*(\@points/10)\relax>0
            taškai%
          \else
            taškų%
          \fi
        \fi
      \fi
    \else
      \ifnum\@points=1
        taškas%
      \else
        \ifnum\numexpr\@points-10*(\@points/10)\relax>0
          taškai%
        \else
          taškų%
        \fi
      \fi
    \fi
  \endgroup
}
\makeatother
\pointname{ \points}
\marginpointname{ \points}
\firstpagefooter{}{\thepage}{}
\runningfooter{}{\thepage}{}
\setlength{\parindent}{0pt}
\firstpageheader{\small\begin{tabular}{@{}l@{\hspace{0.2cm}}l@{}}Mokinio kodas: & \irasolaukas{2.3cm} \\[0.15cm] Klasė: & \irasolaukas{2.3cm}\end{tabular}}{\small\textbf{Atsiskaitymas 1}\\2026.10.28}{Euklido variantas}
\runningheader{\small Kvadratiniai daugianariai ir nelygybės}{}{\small Euklido variantas}
\begin{document}
\begin{center}\textbf{Kvadratiniai daugianariai ir nelygybės}\end{center}
Visuose uždaviniuose laikykite, kad $x\in\mathbb R$. Sprendimus pagrįskite.
\medskip
\begin{questions}
\question[11] Duotas kvadratinis daugianaris $P(x)=2x^2-8x-10$.\droppoints
\begin{tasks}[label=\alph*),label-width=1.2em](1)
\task Nustatykite daugianario šaknis.
\task Užrašykite šį daugianarį dar dviem skirtingais pavidalais.
\task Nustatykite daugianario ekstremumo tašką, ekstremumą ir jo tipą.
\task Nustatykite, su kuriomis $x$ reikšmėmis daugianario reikšmės yra:\\ teigiamos; neigiamos; neteigiamos; neneigiamos.
\end{tasks}
\sprendimas{8}\atsakymas[12]
\vspace{0.25cm}
\question[7] Duotas reiškinys $\displaystyle T(x)=\frac{\sqrt{-x^2+4x+5}}{5-x}$.\droppoints
\begin{tasks}[label=\alph*),label-width=1.2em](1)
\task Apskaičiuokite $T(2)$, $T(5)$, $T(6)$, jeigu šios reikšmės apibrėžtos realiųjų skaičių aibėje.
\task Nustatykite visas realiąsias $x$ reikšmes, kurias galima įstatyti į reiškinį $T(x)$, kad gautume realųjį skaičių.
\end{tasks}
\sprendimas{5}\atsakymas[12]
\newpage
\question[5] Duotas kvadratinis daugianaris $P(x)=2x^2-4x+5$.\droppoints
\par\medskip Dviem skirtingais būdais įrodykite, kad $P(x)>0$ su visomis realiosiomis $x$ reikšmėmis:
\begin{tasks}[label=\alph*),label-width=1.2em](1)
\task išskirdami pilnąjį kvadratą;
\task remdamiesi tik daugianario koeficientais $a$, $b$, $c$ ir diskriminantu.
\end{tasks}
\sprendimas{7}
\vspace{0.35cm}
\question[3] Mokytojas pakvietė prie lentos Jonuką rasti visas $x$ reikšmes, su kuriomis nelygybė yra teisinga:\droppoints
\[(x-3)(x+1)>0.\]
Jonukas sako:\par
\hspace{0.5cm}--- Easy! Padalijame abi puses iš $(x-3)$, gauname $x+1>0$, taigi $x>-1$.\par
Mokytojas sako:\par
\hspace{0.5cm}--- Sėsk, du!\par\medskip
Kur problema Jonuko sprendime? Atsakymą pagrįskite ir išspręskite nelygybę teisingai.
\sprendimas{7}\atsakymas[9]
\newpage
\question[5] Kvadratinis daugianaris $P(x)=ax^2+bx+c$ nagrinėjamas uždarame intervale $x\in[0;9]$. Žinoma, kad $P(1)=P(5)$ ir $b>0$. Nustatykite, su kuria $x$ reikšme šiame intervale $P(x)$ įgyja mažiausią reikšmę ir su kuria -- didžiausią.\droppoints
\sprendimas{8}\atsakymas[10]
\vspace{0.35cm}
\question[4] Sukurkite kvadratinį daugianarį $P(x)$, kuris tenkina abi sąlygas:\droppoints
\begin{tasks}[label=\alph*),label-width=1.2em](1)
\task $P(x)<0$ tada ir tik tada, kai $2<x<6$;
\task $P(0)=24$.
\end{tasks}
Užrašykite gautą daugianarį pavidalu $P(x)=ax^2+bx+c$.
\sprendimas{7}\atsakymas[10]
\end{questions}
\newpage
\begin{center}\large\textbf{Juodraštis}\end{center}
\end{document}
